% =============================================================================
% 009 / 068
% Status: raw
% =============================================================================
% Notes:
%
% Source question:
% 我知道吳語小字和你們構建的吳語書面語是兩回事，書面語和語言的思維，是遠遠比表音字型重要得多。
% 
% 但是，在你們團隊建構的吳語字典和寫作指南尚未出爐之前，大部分的人都是吳語初學者吧？更何況，現在連吳語的發音腔調都正在被官話置換當中。
% 
% 問題來了：那些被你認為官話成分極高的吳語方言（城裡的），這些吳語方言的作品（歌曲或影片） ，值得被標上吳語小字的字幕嗎？
% 
% you know
% =============================================================================

\chapter[我知道吳語小字和你們構建的吳語書面語是兩回事，書面語和語言的思維，是遠遠比表音字型重要得多。 但是，在你們團隊建構的吳語字典和寫作指南尚未出爐之前，大部分的人都是吳語初學者吧？更何況，現在連吳語的發音腔調都正在被官話置換當中。 問題來了：那些被你認為官話成分極高的吳語方言（城裡的），這些吳語方言的作品（歌曲或影片） ，值得被標上吳語小字的字幕嗎？ you know]{我知道吳語小字和你們構建的吳語書面語是兩回事，書面語和語言的思維，是遠遠比表音字型重要得多。 \\[0.35em] 但是，在你們團隊建構的吳語字典和寫作指南尚未出爐之前，大部分的人都是吳語初學者吧？更何況，現在連吳語的發音腔調都正在被官話置換當中。 \\[0.35em] 問題來了：那些被你認為官話成分極高的吳語方言（城裡的），這些吳語方言的作品（歌曲或影片） ，值得被標上吳語小字的字幕嗎？ \\[0.35em] you know}
\qaanswer{
使用吳小字純屬個人自由，只要妳開心可以用來標註任何語言的發音。研發iPhone硬體的人總不能規定使用者說「必須這麼用，不能那麼用」吧？
}

